\documentclass[12pt,a4paper]{extarticle}
\usepackage[left=3cm,right=2cm,top=2cm,bottom=2cm,headheight=18pt]{geometry}
\usepackage{fontspec}
\setmainfont{Liberation Serif}
\setsansfont{Liberation Serif}
\setmonofont{Liberation Mono}
\usepackage{unicode-math}
\setmathfont{texgyretermes-math.otf}
\usepackage{amsmath,graphicx,booktabs,longtable,array,enumitem,pgfplots}
\usepgfplotslibrary{groupplots}
\pgfplotsset{compat=1.17,every axis/.append style={font=\fontsize{11bp}{13.2bp}\selectfont}}
\usepackage[english,russian]{babel}
\usepackage{titlesec,fancyhdr}
\usepackage[hidelinks,unicode]{hyperref}
\urlstyle{same}
\setlength{\parindent}{1.27cm}
\setlength{\parskip}{0pt}
\linespread{1}
\setlength{\emergencystretch}{2em}
\clubpenalty=10000
\widowpenalty=10000
\titleformat{\section}{\normalfont\bfseries}{\thesection.}{0.5em}{}
\titleformat{\subsection}{\normalfont\bfseries}{\thesubsection.}{0.5em}{}
\titlespacing*{\section}{0pt}{1.2em}{0.5em}
\titlespacing*{\subsection}{0pt}{0.9em}{0.3em}
\setlist{nosep,leftmargin=1.27cm}
\pagestyle{fancy}\fancyhf{}
\renewcommand{\headrulewidth}{0pt}
\renewcommand{\footrulewidth}{0pt}
\fancyhead[C]{\ifnum\value{page}>1\thepage\fi}
% Liberation Serif — disclosed substitute, not Times New Roman.
\begin{document}
\begin{center}
\textbf{НАУЧНО-ПРАВОВОЕ ОБОСНОВАНИЕ}\par
Предлагаемые положения о проверке расчётов для акта по Социальному реестру во исполнение PF-193\par
\medskip
Редакция от 1 октября 2026 года; опорная дата расчётов — 30 сентября 2026 года\par
Abduxoliq Akmal ugli Ashuraliyev\par
Самостоятельное исследовательское предложение; принадлежность к учреждению не заявляется\par
\medskip
Основная русская редакция авторского предложения для подачи. Согласованные узбекская и английская редакции приложены; узбекская предназначена для работы над официальным проектом на государственном языке.
\end{center}
\section{Правовое основание, ответственные разработчики и цель}
Настоящий информационно-аналитический материал обосновывает предложение гражданина Abduxoliq Akmal ugli Ashuraliyev для рассмотрения Национальным агентством социальной защиты при Президенте Республики Узбекистан и Министерством экономики и финансов Республики Узбекистан в порученном проекте. Автор не представляет пояснительную записку от имени официальных разработчиков. Цель — обосновать семь основных предлагаемых положений и приложение из двадцати восьми пунктов о проверке расчётов. Основная русская пояснительная записка и её согласованные узбекский и английский переводы организуют сведения для представления по пунктам 83–85 Единой методики. \cite{ref7}.

Пункты 13–14 PF-193 от 11 сентября 2026 года вводят многомерную оценку с 1 января 2027 года и поручают Национальному агентству социальной защиты и Министерству экономики и финансов внести проект в Администрацию Президента до 1 декабря 2026 года. Пункт 14 не определяет окончательный вид акта. Поэтому предлагаются положения для включения и согласованные изменения, которые интегрирует компетентный разработчик; не утверждается, что отдельное постановление Кабинета Министров само по себе полностью исполняет поручение. \cite{ref1}.

Классификация семьи включает как минимум три разных действия: вычисление балла или дохода; присвоение категории по правовым правилам; определение права на конкретную программу и размера выплаты. Если объединить их в один непрозрачный результат, трудно установить, возник ли спор из-за исходных данных, арифметики, исключения, категории или отдельного условия программы. Воспроизводимая система сохраняет это разграничение, даже если все этапы выполняет одна информационная система.

Постановление Кабинета Министров № 35 от 29 января 2026 года регулирует ведение реестра, доход, повторную оценку, обмен информацией и обжалование. Проверенная на рабочую дату консолидированная редакция включает изменения 2026 года. PF-258 от 26 декабря 2025 года содержит вышестоящие положения о категориях и тяжёлых жизненных обстоятельствах. Они существенны до их законного изменения. Новая модель сама по себе не отменяет оснований отказа и не даёт программисту полномочия устранять противоречия правового текста. \cite{ref2}; \cite{ref3}.

Пункт 13 ПФ-193 подтверждает измерения оценки: финансовое положение, хронические заболевания, инвалидность, внешнюю трудовую миграцию и многодетность; распределение по категориям на основе баллов предусмотрено с 1 января 2027 года. Пункты 13–14 этого указа не устанавливают численные веса, преобразования, округление и пороги категорий, а поручают подготовить соответствующий акт. Указ, ведомственный перечень обсуждаемых проектов и открытый индексированный поиск проверены 1 октября 2026 года; полная численная спецификация 2027 года не получена. Это результат указанного поиска, а не доказательство отсутствия внутренней либо опубликованной спецификации. В обращении по статье 22 запрашиваются датированная спецификация и её утверждение либо статус проекта. \cite{ref1}.

Исследовательский вопрос допускает опровержение: \textbf{дают ли выбранные опубликованные правила воспроизводимые и внутренне согласованные условия расчёта дохода и категории на заданных границах и в объявленных областях входных значений?} Ошибка переноса правила, пропущенное применимое исключение, более узкая законная область или расхождение с независимой арифметикой опровергнут соответствующий вывод. Вопрос не касается точности новой обученной модели бедности: веса 2027 года и записи Агентства не предоставлены.

Статья 22 Закона ЗРУ-682 допускает участие согласных граждан с необходимыми знаниями и опытом в рабочих группах. Назначение зависит от решения разработчика; квалификацию следует обосновать. Статья 20 позволяет представить предложение. Статья 24 регулирует общественное обсуждение и обоснование отклонения предложений. \cite{ref7}.

\section{Действующее регулирование и относящиеся к вопросу доказательства}
\subsection{Получение источников и актуальность}
Это целевой повествовательный обзор первичных документов, а не систематический обзор. В проекте получены официальные консолидированные акты, статистика, страницы издателей, опубликованные исследования с авторских сайтов и отчёты Всемирного банка. Изучены соответствующие положения, аннотации и выбранные разделы полных текстов. Источники, доступные только в поисковом фрагменте, исключены из содержательного обобщения. Прямое получение использовано, поскольку провайдер research-lookup не был настроен. Поиск последующих исправлений и действующих актов проверяет актуальность, но не доказывает отсутствие неопубликованного проекта или неиндексированного изменения.

Правовые тексты указаны по акту и положению. Для крупных отчётов различаются печатная и электронная нумерация: печатная страница 40 оценки Узбекистана 2026 года соответствует странице 42 PDF, страница 141 справочника систем социальной защиты — странице 171 PDF. Источники обосновывают архитектуру проверки, но не устанавливают её единственную оптимальность. Иностранные исследования помогают определить вопросы и механизмы; их коэффициенты и ожидаемые эффекты не переносятся на Узбекистан.

\subsection{Современные данные об Узбекистане}
SIAT публикует национальные уровни бедности 17,0; 14,1; 11,0; 8,9 и 5,8 процента за 2021–2025 годы. Это контекстные агрегаты, а не данные калибровки аудита или доказательство того, что конкретный алгоритм реестра вызвал снижение. Региональные уровни не проверяют правильность индивидуальной классификации. \cite{ref6}.

Опубликованная 18 сентября 2026 года оценка бедности и равенства возможностей Всемирного банка обсуждает охват и адресность помощи малообеспеченным семьям. На печатной странице 40 отмечено отсутствие на тот момент доказательств эффективности формулы, пересмотренной в начале 2026 года. Техническое приложение описывает статистический показатель благосостояния на основе дохода. Он не тождествен административному показателю с нормативно вменённым доходом и специальными включениями и исключениями. Отчёт подтверждает потребность в эмпирическом исследовании, но не проверяет спецификацию баллов 2027 года, не полученную для данного проекта и не устанавливает операционную частоту наших контрпримеров. \cite{ref11}.

\subsection{Наблюдаемые сведения об услуге и пределы вывода}
Сообщение Правительства от 17 июля 2026 года указывает 2\,000 необоснованных отказов в 184 махаллях за месяц и рост обращений Президенту по реестру с 30\,000 до 46\,000; поручен прокурорский надзор. Это основание изучать причины решения, но не атрибуция программе, национальный знаменатель, текущая частота или эффект предлагаемого контроля. Причину нельзя вывести из формульного контрпримера. \cite{ref28}.

Официальная страница услуги 587, просмотренная 1 октября 2026 года, описывает канал заявления/результата, бесплатность, отсутствие требуемых документов и указанный месячный срок. Показано 640\,876 общих статусов, включая 16\,564 в процессе. Не приведены период накопления, возраст ожидающих дел, время обработки, дата выгрузки или доступность работников. Это не годовой спрос, число просрочек или мощность. Сообщение Ургенча от 19 февраля документирует местное практическое обучение по постановлению № 35, законному включению/исключению и надёжности данных; общенациональная кадровая готовность не установлена. Могут использоваться действующие каналы и обучение; их дополнительная нагрузка требует реальных записей. \cite{ref29}; \cite{ref30}.

Записи выпусков, время обработки, покрытие договорами и обеспеченные средствами задания Агентства не получены. Публичные страницы финансовой прозрачности и более широкие социальные проекты не установили выделение средств или свободную мощность этого модуля. Поэтому нет прогноза очереди, оценки штата или гарантии срока исполнения. Если интеграция меняет обработку заявлений, Агентство использует датированные поступления/завершения, возраст ожидающих дел, задержки источников и время обработки для разрешённой проверки мощности до обязательства изменения. Общий бюджет или штатные места не заменяют наблюдения.

\subsection{Значение международных исследований}
Обзор социальных реестров Всемирного банка 2025 года рассматривает динамическое обновление, доступ по обращению, совместимость систем и административные возможности. Реестр является частью предоставления поддержки; охват базы и право на программу различны. Переносимого порога приёмки для Узбекистана обзор не даёт. \cite{ref12}.

Оценка качества данных в Индонезии подчёркивает единые стандарты, актуальность, аудит и местные возможности исполнения. Это основания для паспорта переменной и ясного пути исправления. Возможный переносимый механизм — прослеживаемость исходных значений; величина пользы зависит от систем и кадров Узбекистана, которые не наблюдались. \cite{ref13}.

Справочник систем предоставления социальной защиты разграничивает оценку, право, включение, уведомление и выплату. Рассмотрение жалоб охватывает ошибки разных этапов. Поэтому следует сохранять расчёт и направлять исправление в ответственный процесс, а не приравнивать балл к выплате. Исторические примеры стран не представлены как их нынешнее устройство. \cite{ref14}.

Обобщение исследований адресности 2022 года выстраивает иерархию целей политики, устройства программы, выбора адресности и исполнения. Польза усложнения моделей зависит от данных и контекста. Это основание сначала определить смысл правильности, затем оптимизировать балл. \cite{ref15}.

Два первичных исследования показывают необходимость явного критерия благосостояния. Рандомизированное исследование Alatas и соавторов сравнивало косвенную проверку дохода, оценку сообществом и смешанный подход в 640 индонезийских деревнях: ранжирование по потреблению и удовлетворённость участников не определяли один и тот же предпочтительный метод. Это результат одного контекста, не оценка Узбекистана. Исследование Aiken и соавторов в Того показало зависимость результатов адресности по телефонным данным от сравнения с географической адресностью либо гипотетическим всеобъемлющим реестром. Нельзя считать старое административное решение безусловной истинной меткой нуждаемости. \cite{ref16}; \cite{ref17}.

Различаются \textbf{соответствие правилу}, \textbf{качество данных}, \textbf{достоверность оценки благосостояния} и \textbf{результаты предоставления помощи}. Выполненный анализ устанавливает отдельные свойства соответствия. Остальные направления — будущая программа Агентства.

\section{Исходные данные и метод обосновывающего расчёта}
\subsection{Объекты данных и единицы}
Входы — опубликованные правовые константы и сконструированные случаи. Месячные минимальные потребительские расходы на человека в 2026 году обозначены \(P=715\,000\) сумов; минимальный размер оплаты труда с 1 сентября 2026 года — \(M=1\,360\,000\). Первое значение взято из сообщения Национального комитета по статистике от 4 марта 2026 года, второе — из PF-115. Значение \(M\) используется только в датированной диагностике одного сентябрьского месяца и не подставляется задним числом во все месяцы трёхмесячной истории. \cite{ref5}; \cite{ref4}.

Пусть \(Y\) — учитываемый доход семьи за выбранные три месяца, \(N\) — юридически учитываемое число членов. Месячный доход на человека \(D=Y/(3N)\). Полная агрегация \(Y\) из административных данных здесь не восстановлена: проверяются выбранные формулы вкладов и условия категорий. Месячные вклады один раз переводятся в сумму периода до деления на три; повторное деление или пропуск преобразования меняет смысл \(D\).

Эталонная реализация — \nolinkurl{scripts/run_analysis.py}. Числа задаются целыми или строками рациональных и десятичных чисел; расчёт использует точный тип Python Fraction. Отклоняются числа с плавающей точкой, неверные делители, неправильные значения, повторяющиеся поля и неподдерживаемые поля конфигурации. Точная арифметика исключает случайное двоичное приближение в проверке равенства. Производственное округление всё равно требует утверждённого правила; соглашение Агентства здесь не выдумано.

\subsection{Предварительный план и выбор случаев}
Первоначальный аудит задан в предоставленном предварительно заданном плане до исполнения после ознакомления с нормами. Кандидаты были известны заранее; это формальная диагностика, не слепое подтверждающее исследование населения. Расширение сохранило утверждённый метод, случаи зафиксированы до исполнения. Контрпример первого независимого рецензента обозначен как полученный после первоначального анализа и сохранён как регрессионная проверка.

Первоначальная сетка дохода включает \(P\), \(1{,}5P\), \(2P\) со сдвигами минус один, ноль и плюс один сум и контроль прежней категории при \(P\). Расширение проверяет тяжёлые обстоятельства около \(P\), \(4P/3\), \(1{,}5P\), \(2P\); все целые площади построек \(B\) от нуля до \(U\) при \(U=100,1000,1001,2000\) м²; неизвестные переводы и значения около \(2M\) для двух групп стран при двух условных периодах. Налог проверяется около алгебраической точки перехода со сдвигом в один сум во всех восьми сочетаниях территории и пола. Неформальный доход также охватывает восемь сочетаний. Случаи коз: 0, 1, 14, 15, 20 и 100.

У сеток нет выборочных весов и генеральной совокупности. Число нарушений — диагностика программы и области. Доверительный интервал распространённости среди семей Узбекистана был бы бессмысленным: семьи не отбирались. Перебор исчерпывающий только внутри заявленной сконструированной области, а не для всех правовых и операционных состояний.

\subsection{Правовые вычислительные условия}
Для нового трудоспособного заявителя без других оснований отказа, прежней основной категории и рекомендации по тяжёлой ситуации пункт 33(a) отказывает в первичном признании по доходу при \(D>P\). Пункт 32 предусматривает признание при отсутствии оснований пункта 33, но пункт 39(b) описывает доход основной категории как \(D<P\). Условия проверены отдельно; полное решение по заявлению не закодировано.

Для тяжёлой ситуации \(A=3D/4\). При соответствующих обстоятельствах, действительной рекомендации, наличии квоты, повторной оценке и отсутствии иных применимых оснований пункты 34–35 дают доходный путь \(D>P\) и \(A\leq1{,}5P\). Общий пункт 39 задаёт интервал от \(P\) до \(1{,}5P\), не раскрывая связь с уменьшением дохода. Проверены прочтение через обычный и через уменьшенный доход. Применение приоритета специальной нормы может разрешить вопрос; расчёт показывает, какую область такой приоритет должен сохранять.

Буквальная производная площадь земли \(L=U-B-\min\{U/5,200\}\) в м²; заявленная алгебраическая область \(0\leq B\leq U\). Месячный нормативный доход \(L\cdot0{,}035M\cdot k/100\), где 100 м² составляют одну сотку, \(k\) — соответствующий местный налоговый коэффициент. Диагностическое допущение \(k=1\) явно объявлено: фактические коэффициенты махаллей недоступны. Предлагаемый предел \(\max(0,L)\) рассчитан отдельно от буквального правила.

По пункту 26 неизвестные либо меньшие \(2M\) переводы приводят к вменению месячного дохода \(cM\): \(c=3\) для перечисленных стран и \(c=2\) для других. Период сравнения не установлен в полученной спецификации. Условная модель A считает сравниваемую величину \(q\) месячной; B — суммой трёх месяцев, переводя принятые фактические переводы в \(q/3\) за месяц. Обе сохраняют предусмотренный месячный нормативный вклад. Это демонстрация необходимости определить единицу времени; ни одна модель не объявлена действующим алгоритмом Агентства.

По пункту 24 доход предпринимателя \(E=\max\{5T/6,B_0\}\), где \(T\) — налог за три месяца, \(B_0\) — месячная база пункта 23. Пункт 25 задаёт \(1{,}05B_0\) при выполнении условий неформальной занятости. По пункту 28 коза равна 0,07 условной головы крупного рогатого скота; первая условная голова освобождается. Формула только для коз: \(G=\max\{0{,}07g-1,0\}M/2\). Другие коэффициенты животных из постановления № 689 не перенесены и не исполнены.

\subsection{Проверка результатов}
Проверка объединяет официальный пример, ручные неравенства, независимую целочисленную арифметику и исполняемые регрессионные тесты. \nolinkurl{scripts/verify_results.py} не импортирует расчётный движок: проверяет все 4\,105 земельных, 12 случаев тяжёлой ситуации, 24 налоговых, 16 условных переводов и 6 случаев коз эквивалентными кусочными либо целочисленными выражениями. Научных тестов 27; ещё 10 проверок происхождения и независимого контроля отклоняют устаревший, изменённый или неполный результат. Всего 37: равенство, обе стороны границ, неизвестное и ноль, отрицательные веса, неправильный ввод, формулы, количества случаев и условные интерпретации. Первый рецензент выявил два дефекта раннего предложения: общий нижний предел исключал допустимый случай тяжёлой ситуации и было пропущено положение приложения к указу. Оба исправлены, контрпример сохранён.

Вычислительные материалы прошли независимую внутреннюю рецензию и исполняемую проверку. Текущий набор содержит 37 регрессионных тестов; отдельная арифметическая проверка контролирует зафиксированный состав случаев. Рецензия подтверждает изложенные вычислительные результаты в указанном объёме. Она не устанавливает, что производственная система Агентства исполняет те же правила или что не полученная для этого проекта модель 2027 года проверена.

\section{Результаты выполненного аудита правил}
\subsection{Официальный пример и равенство}
Опубликованный пример пункта 30 даёт \(6\,480\,000/(5\cdot3)=\mathbf{432\,000}\) сумов на человека в месяц; отличие от опубликованного ответа нулевое. Это подтверждает арифметику примера. Доход за май–июль относится к заявлению сентября 2025 года. Точный календарный сдвиг следует сохранить в спецификации, а не угадывать из общей фразы.

При \(D=P=715\,000\) оба высказывания \(D>P\) и \(D<P\) ложны. В сконструированных условиях первичное признание пункта 32 не отклоняется по доходу, а пункт 39(b) не даёт соответствующей категории. Несогласование относится к равенству, не всей окрестности. Контроль прежней основной категории при \(P\) удовлетворяет обычному пограничному интервалу, что показывает необходимость учитывать прежний статус.

Пункт 3 PF-258 также применяет строгое условие ниже \(P\) для основных категорий. Поэтому замена «ниже» на «не превышает» только в постановлении Кабинета Министров не согласует иерархию. Вариант включения равенства требует согласованного изменения пункта 3(a)–(b), примечания (a) приложения 1 к PF-258 и пункта 39(a)–(b) постановления № 35, а также обычной пограничной ветви с сохранением специальной ветви тяжёлой ситуации. Аудит не выбирает обязательную политику равенства; выбранная политика должна быть явной и согласованной.

\subsection{Интервалы тяжёлой ситуации: доказательство и контрпример}
Условия специального пути упрощаются точно:
\begin{equation}
D>P,\quad\frac34D\leq\frac32P\quad\Longleftrightarrow\quad P<D\leq2P.
\end{equation}
При обычном доходе пересечение общего интервала со специальным равно \(P<D\leq1{,}5P\); исключается \(1{,}5P<D\leq2P\). При уменьшенном доходе нижний предел \(A\geq P\) равнозначен \(D\geq4P/3\), исключая \(P<D<4P/3\). Ни подстановка \(D\), ни подстановка \(A\) в один общий интервал не воспроизводит специальный путь целиком. Машинное описание требует исключения или приоритета ветви. Это вопрос совместного толкования с разработчиком, не доказательство фактических отказов этим семьям.

При \(P=715\,000\) границы равны \(4P/3=2\,860\,000/3\), \(1{,}5P=1\,072\,500\), \(2P=1\,430\,000\) сумов на человека в месяц. Проверены 12 сконструированных значений; точные рациональные числа сохраняются, даже если график округляет координату.

Контрпример рецензента: три члена и доход 8\,580\,000 сумов за три месяца. Обычный \(D=2\,860\,000/3\), уменьшенный \(A=715\,000\). При всех законных условиях тяжёлой ситуации верхняя граница пункта 35(b) соблюдается. Отвергнутый вариант с уменьшенным доходом строго выше \(P\) исключил бы случай. Исправленный текст разделяет обычную ветвь прежней категории и тяжёлую ситуацию и не добавляет выдуманный нижний предел уменьшенного дохода.

\begin{center}\begin{tikzpicture}
\begin{axis}[width=.94\linewidth,height=6cm,xmin=.85,xmax=2.1,ymin=.5,ymax=3.65,xtick={1,1.333333333,1.5,2},xticklabels={1,$4/3$,$3/2$,2},ytick=\empty,axis y line=none,axis x line=bottom,xlabel={Обычный месячный доход на человека / $P$}]
\addplot[black,line width=3pt] coordinates {(1,3)(2,3)};
\addplot[black,line width=3pt,dashed] coordinates {(1,2)(1.5,2)};
\addplot[black,line width=3pt,dotted] coordinates {(1.333333333,1)(2,1)};
\addplot[only marks,mark=*,mark size=3pt] coordinates {(2,3)(1.5,2)(1.333333333,1)(2,1)};
\addplot[only marks,mark=o,mark size=3pt,mark options={fill=white}] coordinates {(1,3)(1,2)};
\node[anchor=west] at (axis cs:.88,3.28) {Специальный доходный путь};
\node[anchor=west] at (axis cs:.88,2.28) {Пересечение: обычный доход};
\node[anchor=west] at (axis cs:.88,1.28) {Пересечение: уменьшенный доход};
\end{axis}\end{tikzpicture}\end{center}
\textbf{Рисунок 1.} Доход нормирован на \(P\). Закрашенная граница включена, пустая исключена. Специальный путь дополнительно требует соответствующих обстоятельств, рекомендации, квоты и иных условий. Пересечения показывают различие толкований, а не наблюдаемые отказы.

\subsection{Аудит области земельной формулы}
При \(U\leq1000\) формула \(L=4U/5-B\), поэтому \(L<0\) точно при \(B>4U/5\). При \(U>1000\): \(L=U-B-200\), отрицательность точно при \(B>U-200\). Неравенства независимо проверяют перебор. В области \(0\leq B\leq U\) нижний предел \(-\min(U/5,200)\). Физическая площадь не может быть отрицательной, но операционные кадастровые ограничения могут исключать сконструированные входы.

\begingroup\fontsize{11bp}{13.2bp}\selectfont
\begin{longtable}{p{\dimexpr\linewidth/5-2\tabcolsep\relax} p{\dimexpr\linewidth/5-2\tabcolsep\relax} p{\dimexpr\linewidth/5-2\tabcolsep\relax} p{\dimexpr\linewidth/5-2\tabcolsep\relax} p{\dimexpr\linewidth/5-2\tabcolsep\relax}}
\toprule
Участок $U$ (м²) & Целые случаи & Отрицательные & Минимум $L$ (м²) & Максимум изменения (сум/месяц) \\
\midrule
\endhead
100 & 101 & 20 & -20 & 9,520 \\
1,000 & 1,001 & 200 & -200 & 95,200 \\
1,001 & 1,002 & 200 & -200 & 95,200 \\
2,000 & 2,001 & 200 & -200 & 95,200 \\
Вся сетка & 4,105 & 620 & — & — \\
\bottomrule
\end{longtable}
\endgroup

\textbf{Таблица 1.} Расчёт при сентябрьском \(M=1\,360\,000\) и диагностическом \(k=1\). Количества относятся к сконструированной сетке. Последний столбец — изменение нормативного месячного дохода семьи, не сокращение выплаты и не бюджетная оценка.

При \(U=100\), \(B=90\) получено \(L=-10\) м², буквальный месячный доход минус 4\,760 сумов; предлагаемый предел даёт ноль. В заявленной области предел увеличивает месячный вклад максимум на \(95\,200k\) при сентябрьском \(M\); деление на \(N\) даёт вклад на человека. Исправление отрицательного вклада может повысить доход и не является автоматическим увеличением помощи. Последствия для категории и выплаты требуют законного повторения до и после изменения и правил соответствующей программы.

Нулевой предел — обоснованное математическое предложение, но не замена определения геометрии. Кадастровый источник должен уточнить: является ли \(B\) площадью застройки, какие объекты входят, пересекаются ли вычитаемые площади, совпадают ли участок и дата для \(U\) и \(B\). Если 20-процентный вычет соответствует определённым видам использования, непересекающееся геометрическое построение может быть предпочтительнее. Значение определяет разработчик; кадастровые факты не выдуманы.

\begin{center}\begin{tikzpicture}
\begin{groupplot}[group style={group size=2 by 2,horizontal sep=1.4cm,vertical sep=1.8cm},width=.44\linewidth,height=5cm,xlabel={$B$ (m²)},ylabel={$L$ (m²)},scaled ticks=false]
\nextgroupplot[title={$U=100$; 20 отрицательных случаев},xmin=0,xmax=100]
\addplot[black,thick] coordinates {(0,80)(100,-20)};
\addplot[black,dashed,thick] coordinates {(0,80)(80,0)(100,0)};
\nextgroupplot[title={$U=1000$; 200 отрицательных случаев},xmin=0,xmax=1000]
\addplot[black,thick] coordinates {(0,800)(1000,-200)};
\addplot[black,dashed,thick] coordinates {(0,800)(800,0)(1000,0)};
\nextgroupplot[title={$U=1001$; 200 отрицательных случаев},xmin=0,xmax=1001]
\addplot[black,thick] coordinates {(0,801)(1001,-200)};
\addplot[black,dashed,thick] coordinates {(0,801)(801,0)(1001,0)};
\nextgroupplot[title={$U=2000$; 200 отрицательных случаев},xmin=0,xmax=2000]
\addplot[black,thick] coordinates {(0,1800)(2000,-200)};
\addplot[black,dashed,thick] coordinates {(0,1800)(1800,0)(2000,0)};
\end{groupplot}\end{tikzpicture}\end{center}
\textbf{Рисунок 2.} Каждая панель относится к заранее заданному \(U\) и всем целым \(B\) от нуля до \(U\); вертикальная ось — м². Отрицательная область следует буквальной формуле, штриховая линия — предлагаемый предел. График не добавляет административных наблюдений.

\subsection{Единицы переводов и условные скачки}
Равенство \(q=2M\) не удовлетворяет условию вменения «ниже порога». При месячном прочтении вклад перечисленных стран меняется с \(3M\) непосредственно ниже порога на \(2M\) при равенстве: месячное снижение \(M=1\,360\,000\). Для прочих стран переход с \(2M\) на \(2M\) непрерывен. При трёхмесячной сумме фактический месячный доход при равенстве \(2M/3\), что даёт другие скачки. Шестнадцать случаев сохраняют обе трактовки, обе группы и неизвестное значение.

\begingroup\fontsize{11bp}{13.2bp}\selectfont
\begin{longtable}{p{\dimexpr\linewidth/4-2\tabcolsep\relax} p{\dimexpr\linewidth/4-2\tabcolsep\relax} p{\dimexpr\linewidth/4-2\tabcolsep\relax} p{\dimexpr\linewidth/4-2\tabcolsep\relax}}
\toprule
Условный период & Группа стран & Снижение за месяц (сум) & На человека при четырёх членах (сум) \\
\midrule
\endhead
Месяц & Перечисленные & 1,360,000 & 340,000 \\
Месяц & Прочие & 0 & 0 \\
Три месяца & Перечисленные & 9,520,000/3 & 2,380,000/3 \\
Три месяца & Прочие & 5,440,000/3 & 1,360,000/3 \\
\bottomrule
\end{longtable}
\endgroup

\textbf{Таблица 2.} Точная условная арифметика в сумах. Период сравнения не разрешён; строки не являются результатом проверки официальной реализации. Изменяется доходный вклад, не размер помощи.

Месячное правило перечисленных стран немонотонно при переходе от нормативного к фактическому доходу: более высокий сообщённый перевод может дать меньший учитываемый вклад. Это не обязательно ошибка нормы: резервная оценка и наблюдение могут намеренно различаться. Однако нужны научное и правовое объяснение, воспроизводимое применение источника и оценка изменений категории. Программист не должен молча заменять правило на максимум фактического и нормативного дохода.

Неизвестное и записанный ноль различаются, даже если действующая ветвь даёт одинаковое вменение: различны причины, доказательства и исправление. Сохраняются период записи, объединение отправителей, дата курса и причина резервного расчёта. Выбор единицы времени или дополнительного минимума ради более гладкого графика — необоснованный выбор политики.

\begin{center}\begin{tikzpicture}
\begin{groupplot}[group style={group size=2 by 1,horizontal sep=1.5cm},width=.44\linewidth,height=6cm,xmin=1,xmax=3,ymin=.5,ymax=3.3,xtick={1,2,3},ytick={1,2,3},xlabel={$q/M$},ylabel={Месячный вклад / $M$}]
\nextgroupplot[title={Месячное сравнение}]
\addplot[black,thick] coordinates {(1,3)(2,3)};
\addplot[black,dashed,thick] coordinates {(1,2)(2,2)};
\addplot[black,dashed] coordinates {(2,.5)(2,3.3)};
\addplot[black,thick] coordinates {(2,2)(3,3)};
\addplot[only marks,mark=*,mark size=3pt] coordinates {(2,2)};
\addplot[only marks,mark=o,mark size=3pt,mark options={fill=white}] coordinates {(2,3)};
\nextgroupplot[title={Сумма трёх месяцев}]
\addplot[black,thick] coordinates {(1,3)(2,3)};
\addplot[black,dashed,thick] coordinates {(1,2)(2,2)};
\addplot[black,dashed] coordinates {(2,.5)(2,3.3)};
\addplot[black,thick] coordinates {(2,.666666667)(3,1)};
\addplot[only marks,mark=*,mark size=3pt] coordinates {(2,.666666667)};
\addplot[only marks,mark=o,mark size=3pt,mark options={fill=white}] coordinates {(2,3)(2,2)};
\end{groupplot}\end{tikzpicture}\end{center}
\textbf{Рисунок 3.} По горизонтали \(q/M\), по вертикали месячный вклад относительно \(M\). Сплошные горизонтальные участки — перечисленные страны, штриховые — прочие; фактическая ветвь совпадает. Панели сохраняют два периода. Вертикальный пунктир — порог. Неизвестное значение сохранено в JSON аудита и не изображено числовой точкой.

\subsection{Коэффициенты, налоговый минимум и дробное поголовье}
Пункт 23 содержит восемь коэффициентов месячной самозанятости: четыре территории и два пола. При сентябрьском \(M\) разница мужского и женского коэффициента в каждой территории \(0{,}4M=544\,000\) сумов. Для сконструированной семьи из четырёх человек один такой член меняет доход на человека на 136\,000. Это чувствительность нормативного дохода, не измерение заработка, дискриминации, неравного права или ошибки. Эмпирическое обоснование и правовые вопросы равенства требуют отдельных доказательств и экспертизы.

\begingroup\fontsize{11bp}{13.2bp}\selectfont
\begin{longtable}{p{\dimexpr\linewidth/5-2\tabcolsep\relax} p{\dimexpr\linewidth/5-2\tabcolsep\relax} p{\dimexpr\linewidth/5-2\tabcolsep\relax} p{\dimexpr\linewidth/5-2\tabcolsep\relax} p{\dimexpr\linewidth/5-2\tabcolsep\relax}}
\toprule
Территория & База: мужчины & База: женщины & Неформальный: мужчины & Неформальный: женщины \\
\midrule
\endhead
Ташкент & 3,672,000 & 3,128,000 & 3,855,600 & 3,284,400 \\
Нукус / областные центры & 3,264,000 & 2,720,000 & 3,427,200 & 2,856,000 \\
Прочие города & 2,856,000 & 2,312,000 & 2,998,800 & 2,427,600 \\
Прочие территории & 2,448,000 & 1,904,000 & 2,570,400 & 1,999,200 \\
\bottomrule
\end{longtable}
\endgroup

\textbf{Таблица 3.} Прямые вычисленные вклады при \(M=1\,360\,000\). Неформальный вклад умножает базу на 1,05 только при условиях пункта 25. Реальный состав семей не выводится.

По пункту 24 налоговое выражение равно базе при \(T=3B_0/2{,}5=6B_0/5\). Ниже результат \(B_0\), выше — рост с наклоном \(5/6\) по отношению к трёхмесячному налогу. Все 24 случая совпадают с кусочным выражением. Минимум непрерывен и монотонен, в отличие от переключения переводов. Одного общего показателя чувствительности к порогу недостаточно для разных правил.

Для 14 коз результат ноль: \(14\cdot0{,}07=0{,}98\) условной головы. Для 15: 1,05 головы, после освобождения 0,05, месячный вклад 34\,000 сумов. Для 20 — 272\,000, для 100 — 4\,080\,000. Округление условных голов вниз перед освобождением изменило бы результаты и здесь не закодировано. Официальное правило округления должно определить дробные эквиваленты. Другие виды и смешанные стада не реализованы.

\subsection{Вопросы спецификации без произвольного числового ответа}
Нужно согласовать включение месяцев и календарный сдвиг пункта 18 с примером пункта 30. Также общие исключения субсидий пункта 20 и прямое включение субсидии продажи возобновляемой электроэнергии \(QP\) в пункте 29 требуют записи приоритета. Специальное включение может правомерно определять применение общего исключения. Не утверждается незаконность \(QP\) либо фактический двойной учёт. Ни один вопрос не превращён в произвольное программное допущение.

\section{Предлагаемые положения и их правовое обоснование}
\subsection{Вид акта и иерархия}
Нормативный проект в материале 1 содержит семь основных пунктов и приложение из двадцати восьми пунктов. Центральная обязанность — проверять соответствие расчёта утверждённой методике, сохранять основания и показывать неразрешённые случаи без нового основания отказа. Добавлены конкретный субъект технической приёмки, её условия, даты применения и законный порядок хранения. Модуль дополняет существующую систему. Окончательный вид, бюджетные решения и согласованные изменения определяет разработчик.

Вариант равенства требует согласованного изменения президентского и правительственного актов; предел площади — изменения формулы постановления № 35. Остальные положения задают проверку и прослеживаемость. Их требования принятия различаются; нельзя представить все как простые программные настройки. Применима иерархия статьи 18 ЗРУ-682. \cite{ref7}.

\begingroup\fontsize{11bp}{13.2bp}\selectfont
\begin{longtable}{p{.21\linewidth}p{.12\linewidth}p{.28\linewidth}p{.28\linewidth}}
\toprule
Вопрос & Пункты & Согласование & Вид предложения \\
\midrule\endhead
Равенство и специальный путь & 5, 12, 16 & PF-258: пункт 3 и примечания; № 35: 32–35, 39, 42 & Необязательная поправка категории; обязательное определение законных ветвей \\
Отрицательная площадь & 4, 12 & № 35: пункт 27; кадастр & Необязательная поправка формулы; обязательный источник геометрии \\
Период и переводы & 4, 8 & № 35: 18–19, 21, 26, 29–30 & Единицы, период и агрегация \\
Сбой источника & 8–11, 18 & № 35: 48; обмен & Временная версия и проверка \\
Эталон и неясность & 12–17 & Реестр, оценка и выплаты & Проверка и знаменатели \\
Основания и пересмотр & 18–20 & SMS и обжалование & Доступ и законное направление \\
Исследование и отчёт & 21–24 & Персональные данные; контроль & Внутренний расчёт и раскрытие \\
Приёмка и переход & 25–28 & Методика, даты, выплаты и хранение & Приёмка без правотворческого делегирования \\
\bottomrule\end{longtable}\endgroup
\textbf{Таблица 4.} Связь результата и предлагаемой нормы. Номера относятся к авторскому приложению, не действующему законодательству. Согласование выполняет разработчик; все последующие программы таблица не охватывает.

\subsection{Паспорт методики}
Каждая значимая переменная должна иметь утверждённый паспорт: правовая цель, достоверный источник, связь с человеком и семьёй, единица, период, допустимая область, преобразование, обработка отсутствия, конфликта и устаревания, датированный параметр. Измеренное значение отличается от нормативного вменения. Даты обновления недостаточно: старая запись может описывать правильный период, а только что полученная — неправильный.

Категории выражаются таблицей решений или эквивалентными формальными правилами. Строка указывает применимость, приоритет, условия, выход и правовое основание. Специальная строка сохраняет обстоятельства, рекомендацию, квоту и повторную оценку: формула дохода сама по себе не даёт права на включение. Раздельно фиксируются признание, категория, право на программу и сумма. Неразрешённость — статус законного уточнения, не выдуманная четвёртая категория нуждаемости.

Сохраняются версия модели и программы, снимки источников, периоды, исключения и шаги расчёта. Открытая часть описания публикуется, ограниченная законно доступна уполномоченным проверяющим. Это сочетает воспроизводимость и конфиденциальность, не создавая права внешнего автора на микроданные.

\subsection{Предлагаемые обязательные механизмы}
Агентство с Министерством готовит утверждённую спецификацию, проверяет её до первого применения и существенного изменения и сохраняет протокол. Существенно изменение, способное повлиять на расчёт или категорию, включая сопоставление источника и временное преобразование. Переименование без изменения смысла может быть несущественным. Пункт 25 предлагает техническую приёмку руководителем Агентства или законно уполномоченным заместителем по заключениям методического подразделения, информационной системы и юридической службы. Это распределение ответственности предлагается, а не объявляется действующим полномочием принимать критерии.

Пункт 26 делает приёмку проверяемой: каждому предварительно заданному граничному и исключительному случаю соответствует определённый законный результат; после утверждённого округления нет различий балла или категории для той же утверждённой версии и одинаковых входов; работают качество данных, доступ к основаниям и направление возражения. Неуспех требует исправления и повторения затронутых тестов. Это критерий соответствия, не целевая точность населения. Нельзя вводить произвольные 95 процентов точности нуждаемости без критерия благосостояния и выборки. Приёмка не приостанавливает законные заявления и не меняет установленную дату внедрения.

Приложение требует проверять равенство, исключения, неполные данные и физически неверные производные значения. Запрещает скрывать неразрешённые случаи в успешном знаменателе и выдавать синтетические испытания за доказательство нуждаемости. Требует выяснять затронутые расчёты при нарушении утверждённого правила. Индивидуальные исправления, выплаты и сроки регулируются применимым правом; нового полномочия автоматической приостановки не предложено.

Пункт 48 постановления № 35 разрешает временное изменение алгоритма при нарушенном обмене. Предложение сохраняет полномочие с записью основания, области, начала, временного правила и условия восстановления; соответствующие расчёты различимы для последующей законной проверки. Параллельного чрезвычайного органа не создаётся. Существующий контроль пункта 49 остаётся отправной точкой.

Основания решения включают вменение и исключения; SMS не воспроизводит расчёт. Статьи 53–54 закона об административных процедурах касаются содержания и фактических и правовых оснований административного акта; их применимость и специальные процедуры уточняются для конкретного решения. Пункт 19 предлагает основания вместе с решением в личном электронном кабинете и доступ через центр «Инсон» для не использующего кабинет. Это новая предлагаемая обязанность, не установленная возможность нынешнего интерфейса. \cite{ref10}.

\subsection{Доказательства приёмки, независимые ответы и охват}
Для успешного теста нужен независимо обоснованный ожидаемый ответ, а не только два совпавших вывода. Barr и соавторы рассматривают проблему определения правильного ответа теста и отличают ответ из спецификации от сравнения с другой реализацией. NIST IR 8397 различает тесты из требований, тесты из структуры кода и регрессионные случаи прежних дефектов. Эти источники поддерживают инженерное устройство проверки; они не являются обязательными узбекскими критериями права на помощь. Они не доказывают, что примеры проекта исчерпывают систему Агентства. \cite{ref21}; \cite{ref22}.

В предложении методическое подразделение и юридическая служба устанавливают правовой вывод ожидаемого результата; его проверяет лицо, отличное от ответственного за проверяемую реализацию. Запись указывает норму и редакцию, предпосылки пути, единицы и периоды входов, промежуточные вклады, этап округления и исход. Например, при $D=P$ копирование вывода программы с $\leq P$ скроет расхождение с опубликованным строгим условием категории $<P$. Напротив, общий нижний предел в специальном пути исключит сохранённый контрпример $0.75D=P$. Соответствие кода правилу и необходимость изменения правила — разные вопросы. Неразрешённое толкование направляется для законного определения; совпадение программ либо большинство голосов его не разрешает.

Пункты 12 и 26 теперь требуют реестра связей правил и случаев, а не позволяют выбранным тестам самим определять полноту. Каждая применимая граница, исключение, путь приоритета и достижимое сочетание имеют ожидаемый результат; неприменимый путь исключается с обоснованием. До выполнения фиксируются идентификаторы случаев и входы. После выполнения реестр сопоставляется с выполненными случаями, включая отсутствующие и повторные идентификаторы. Успех на неполном списке не завершает заявленную область проверки. Покрытие кода даёт дополнительное свидетельство, но процент выполненных операторов не устанавливает охват пропущенных правовых правил.

\begingroup\fontsize{11bp}{13.2bp}\selectfont
\begin{longtable}{p{\dimexpr\linewidth/3-2\tabcolsep\relax} p{\dimexpr\linewidth/3-2\tabcolsep\relax} p{\dimexpr\linewidth/3-2\tabcolsep\relax}}
\toprule
Группа правил & Выполненная проверка и область & Дополнительные доказательства приёмки от Агентства\\
\midrule\endhead
Основные условия дохода & 10 искусственных строк; один контроль прежней категории; предпосылки трудоспособного заявителя & Иные отказы, ветви трудоспособности, даты состава семьи и полные приоритеты признания/категории.\\
Трудная ситуация & 12 граничных строк и один отдельно отмеченный последующий контрпример & Условия ситуации, законная рекомендация, квота, переоценка и сочетания с основным/прежним статусом.\\
Земля & 4\,105 целочисленных случаев четырёх площадей; местный коэффициент 1 & Геометрия одного участка из полномочного источника, дробные площади, местный коэффициент по дате, ограничения источника и округление.\\
Налоговый минимум & 24 строки пересечения восьми территориальных/половых ячеек & Смысл и полнота налоговой суммы, совпадение периодов, меняющиеся параметры и исключения источника.\\
Переводы & 16 строк двух условных периодов и двух групп стран & Один законно определённый период сравнения, агрегирование отправителей, повторы, валютное преобразование и состояния источника.\\
Скот & Шесть случаев только с козами и точным дробным преобразованием & Другие виды, смешанное поголовье, освобождения, период и утверждённое округление; здесь не реализованы.\\
Аддитивная демонстрация & Семь искусственных случаев; четыре сопоставимые заданные категории & Официальная формула, область/зависимости, исключения и законные входы. Демонстрация не проверяет неизвестную нелинейную модель.\\
\bottomrule
\end{longtable}
\endgroup
\textbf{Таблица 5.} Выполненная область и неразрешённые производственные предпосылки. Числа относятся к построенным случаям, не семьям. Это реестр доказательств приёмки, не утверждение о выполнении дополнительной работы.

Промежуточные компоненты проверяются и при совпадении итогового балла и категории. Пропущенный вклад и ошибочно добавленный вклад могут взаимно компенсироваться. Контроль сравнивает причины включения/исключения, периоды источников, преобразования и округление, а также итог. Запись различает утверждённую правовую методику, программный выпуск, сопоставление/конфигурацию и датированные параметры, а также даты правового действия и внедрения. Исправленный выпуск при неизменном правовом правиле сохраняет идентификатор правила и получает отдельную программную запись. Входы решения либо законно сохранённый снимок обеспечивают историческое воспроизведение; одной ссылки на меняющуюся базу недостаточно.

Исправленный независимый проверяющий сверяет идентичность и кратность фиксированных случаев и независимо проверяет все десять строк условий дохода. Шесть дополнительных регрессионных тестов отклоняют пустые списки, замену случая повтором без изменения числа строк, изменённое условие равенства и вход, обозначенный как другой фиксированный случай, а также дробные или булевы входы, принудительно преобразованные в идентификаторы, и принимают неизменный отчёт. Полный набор теперь содержит 37 тестов. Исправлен продемонстрированный дефект проверяющего: прежде пустые списки трудной ситуации, налогов, переводов и коз могли достичь отметки успеха. Неизменные численные результаты проходят усиленную проверку. Производственная программа либо выборка семей этим не сертифицируется.

\subsection{Компетенция обжалования: установленный предел совета}
Пункт 26 приложения 4 к постановлению № 35 ограничивает предмет апелляционного совета анкетами, заключениями и актами «махаллинской семёрки» при рассмотрении заявлений. Пункт 31 допускает отказ по неподведомственному вопросу. Нельзя назвать чисто вычислительный спор жалобой в совет и считать компетенцию установленной. Это предел одного конкретного пути, не доказательство отсутствия иных административных и судебных средств. Пункт 33 устанавливает пятнадцатидневное рассмотрение советом. \cite{ref2}.

Статья 24 закона о персональных данных даёт иной путь: решения только на основе автоматизированной обработки допускаются в предусмотренных случаях, а собственник или оператор объясняет процедуру и последствия, даёт возможность возражения и письменно сообщает результат в течение десяти дней. Такое возражение отличается от исправления исходной записи и спора о махаллинском документе. Статья 11 отдельно требует исправления или дополнения по обращению субъекта в течение трёх дней и немедленного исправления установленной недостоверности; этот срок не заменяется десятью днями для автоматического возражения. Пункт 20 направляет автоматическое возражение собственнику или оператору, исправление — достоверному источнику, жалобу совету — в его предметной компетенции. Сохраняются иные законные пути без дополнительного обязательного досудебного этапа. Так программный пункт не изображает расширение компетенции совета. \cite{ref9}.

Разработчик согласует формы и перенаправление различных вопросов. Уведомление указывает получателя, предмет и применимый срок. Смешанное обращение сохраняет первоначальную дату, вопросы направляются отдельно; перенаправление не обнуляет установленный срок. Это предлагаемые рабочие обязанности для использования гарантий, не выдуманное общее право апелляции.

\subsection{Переход, исправление и законное хранение}
Новая политика меняет правило на будущее; программное исправление восстанавливает исполнение правила прежнего решения; исправление источника меняет доказательство за определённый период. Пункт 27 не позволяет автоматически применять более поздний параметр или новый критерий к истории. Повторный расчёт сам по себе не основание взыскания. Возврат или неблагоприятное изменение решения требуют материального и процедурного правового основания. Статьи 59–61 закона об административных процедурах не позволяют техническим приложением обещать безусловное взыскание либо безусловную невзыскиваемость. \cite{ref10}.

Исходный расчёт необходим для проверки, но не оправдывает бессрочное хранение всех персональных записей. Пункт 28 требует основания и периода по категориям записей, доступа, истории исправлений, законного уничтожения или обезличивания. Статья 17 закона о персональных данных применяется совместно с архивными обязанностями. Точный график требует инвентаризации Агентства; выдуманного срока в годах нет. Открытые таблицы проверяются совместно: безобидные агрегаты могут идентифицировать малую группу при связывании. \cite{ref9}.

\subsection{Участие в подготовке по статье 22}
Статья 20 ограничивает инициирование указанных проектов, когда вопрос решается существующими законными механизмами и административными полномочиями. Предложение вносится в уже порученный пунктом 14 PF-193 проект; поручение не обязывает разработчика включить именно эти положения. Каждая задача сопоставляется с действующим полномочием либо дополнительной обязанностью, требующей надлежащего акта. Внутренние проверочные инструкции могут использовать имеющиеся полномочия; критерий большей юридической силы нельзя изменить настройкой программы или актом меньшей силы. Проверочный модуль и необязательные изменения политики решаются отдельно. \cite{ref1}; \cite{ref7}.

Статья 23 требует изучать относящиеся к предмету исследования и обращения граждан, обобщать и использовать предложения и научные рекомендации. Это даёт путь рассмотрения даже при отказе во включении в рабочую группу. Участники действительного общественного или профессионального обсуждения заранее знакомятся с проектом. Предложения обсуждения рекомендательны и подлежат рассмотрению; статья 24 прямо допускает отклонение с обоснованием причин. Воспроизводимые расчёты усиливают доказательную основу рассмотрения, но ноль программных расхождений или положительная рецензия не доказывают обязанности принять авторскую редакцию. Отдельная пояснительная записка указывает доказательства поступления, обязанности ответа и защиту от незаконного отказа в рассмотрении. \cite{ref7}.

В обращении запрашивается включение норм и рассмотрение участия автора в вычислительном направлении рабочей группы. Достоверное описание выполненной работы обосновывает запрос; текущий проект и неконфиденциальная спецификация обсуждаются, защищённые расчёты выполняет Агентство. Назначение, право доступа или принятие не предполагаются.

Вклад по статье 22 включает формальное описание правил, контрольные случаи, выведенные из правовых норм ожидаемые ответы и проект проверочных положений. Проверка расчётов организационно отделяется от проверяемой реализации. Юридическая служба и компетентный разработчик согласуют текст с применимыми юридическими и антикоррупционными процедурами. Эти внутренние проверки являются результатами подготовки по статье 22. \cite{ref7}.

\subsection{Официальное оформление и разделение документов}
Пункты 43–51 единой методики, приложенной к ЗРУ-682, задают A4, левое поле 3 см, правое 2 см, верхнее и нижнее 2 см, основной шрифт 14 пунктов для неведомственных проектов, отступ 1,27 см и одинарный интервал. Номер страницы сверху по центру начинается со второй страницы, пометка проекта справа сверху на первой. Для аналитических приложений пункт 83 задаёт 11–14 пунктов, левое поле 1–3 см, остальные 1–2 см, отступ 0,50–1,27 см, интервал 1–1,2. Рукопись использует 12 пунктов и таблицы не меньше 11; нормативный текст — 14. Это различные виды документов. Вид окончательного акта не установлен, поэтому модуль остаётся предложением для включения. \cite{ref7}.

Ранее принятые указания № 2352 для внесения в Кабинет Министров задают правое поле 1,5 см, интервал 1,2 и Times New Roman; настройки отличаются от действующей методики закона. Здесь исходная основа — закон, а не заявление об идентичности указаний. В местной среде нет Times New Roman; открыто используется совместимая Liberation Serif. Это замена, не точное соответствие Times New Roman. LaTeX применяется по просьбе автора; законодательство называет Word обычным редактором и не предоставляет официального класса LaTeX. Электронный формат подачи согласовывается с разработчиком. \cite{ref18}.

\section{Дальнейшая проверка на официальных правилах и данных Агентства}
\subsection{Предпосылки и данные}
Это протокол последующей разрешённой работы, не выполненный анализ семей. До извлечения фиксируются спецификация, период, группы, вопросы сравнения, поля и условия выпуска. Документируются законное основание, необходимый объём, защита данных о здоровье и научное обезличивание. Замена идентификаторов не равна обезличиванию; исполнение внутри Агентства — мера защиты, не самостоятельное законное основание. \cite{ref9}.

Исполнение предоставленного кода требует отдельного разрешения оператора после проверки состава, версии, прав доступа и исходящих соединений; защищённые входы не передаются внешним сервисам и не попадают в открытые журналы. Это предлагаемые условия безопасного исполнения пункта 22, а не утверждение об уже настроенной среде Агентства. Требования размещения и обработки проверяются по действующей статье \(27^{1}\) закона о персональных данных в редакции ЗРУ-1125 от 26 марта 2026 года. Законная обработка и раскрытие не возникают из одного факта внутреннего запуска. \cite{ref9}.

Минимальные логические поля: внутренний ключ дела, даты заявления и решения, основание состава семьи, версия, исходные значения и периоды, время извлечения, качество, вменение или исключение, эталонный расчёт, записанная категория и код причины. Разрешённые исправления и пересмотры связываются без утраты исходной записи. Чувствительные переменные берутся лишь при необходимости для законно заданного вопроса; идентифицирующий свободный текст для обычного повторения арифметики не нужен.

Связывание проверяет уникальность идентификатора, отношения один-ко-многим, даты состава семьи и дубли переводов. Сравнение итогового дохода недостаточно, если лишнее включение компенсирует пропущенное вычитание. Согласование компонентов предшествует сравнению категории. В защищённых записях сохраняются точный запрос и хеши снимков; нельзя рассчитывать на неизменность живой базы.

\subsection{Четыре разных вопроса оценки}
\textbf{Соответствие правилу}: разрешённая реализация сравнивается с независимо заданным эталоном на полных сравнимых случаях. Различия балла и категории разделены. \(R\) — число расхождений категорий, делённое на число сравнимых определённых случаев. При нулевом знаменателе показатель не определён. Ноль подтверждает согласие в этом наборе, не правильность оценки нуждаемости.

\textbf{Устойчивость к качеству входа}: заранее заданные исправления или интервалы неизвестных, устаревших и конфликтующих значений меняют выводы. Сценарий указывает поля, допустимые значения, состав случаев и зависимости. Публикуются числитель смен категории, знаменатель парных определённых случаев и исключённые или неразрешённые количества. Матрица переходов информативнее единой доли: направление и исключения важны. Переход после правильного исправления может быть предусмотрен; ненулевая чувствительность сама по себе не дефект.

\textbf{Достоверность нуждаемости}: нужны независимый обоснованный критерий благосостояния и выборка целевого населения. Старые решения и результаты жалоб не являются автоматически истинными метками. Жалобы зависят от доступа, знания и возможности оспаривания; у не подавших жалоб тоже возможны ошибки. Последующее исследование ошибок включения и исключения указывает цель, основу выборки, веса, период измерения, размер и неопределённость. Такое исследование не выполнено.

\textbf{Предоставление помощи}: уведомление, срок решения, нагрузка исправления, обращаемость и фактическое получение. Их нельзя вывести из точности категорий. Изменение времени или затрат требует сопоставимых базовых и последующих наблюдений. Измеренных сбережений, сроков или причинного воздействия проверки на бедность рукопись не заявляет.

OECD/JRC рекомендует проверять зависимость составного показателя от выбора переменных, нормализации, вменения, весов и агрегирования. Руководство относится к составным показателям сопоставления стран; применение к категориям семей здесь является методической адаптацией. Наши заранее заданные сценарии и внешние границы не объявляются выполненным дисперсионным анализом или доказательством допустимости любого веса. \cite{ref23}.

\subsection{Ограниченная неопределённость без выдуманного вменения}
Только для утверждённого аддитивного балла \(S=\sum_jw_jz_j\) с обоснованными \(a_j\leq z_j\leq b_j\):
\[
L=\sum_j\min(w_ja_j,w_jb_j),\qquad U=\sum_j\max(w_ja_j,w_jb_j).
\]
Если переменные независимо принимают любые значения интервалов, границы достижимы. Для дискретных или связанных входов это консервативная внешняя граница, включающая возможные недостижимые категории. Для утверждения достижимости всех категорий нужна проверка зависимостей.

Синтетическая демонстрация реализует лишь этот случай с искусственными переменными и категориями: семь случаев, пять полных и два неразрешённых; четыре сравнимы с заданным эталоном, расхождений в них нет. Это программные примеры с установленными ответами, не измеренная точность узбекской модели. Название «балл» не позволяет выполнять нелинейное дерево или полный процесс реестра.

Точные дроби тоже не решают неопределённость в общем: они определяют арифметику известных входов, но не актуальность записи, правовой приоритет ветви и соответствие показателя тяжёлой ситуации. Неясность источника и нормы сохраняется, а не превращается в внешне точное число.

\subsection{Выпуск и устранение дефектов}
Различается источник несогласия: данные, период, правовая спецификация, округление, код либо намеренная политика. Исправление направляется ответственному процессу. Программный дефект требует новой версии и повторения затронутых испытаний, изменение нормы — компетентного утверждения. Сохраняются обе версии для воспроизведения причины.

Внутренний протокол фиксирует область, версию, случаи, сравнимые и неразрешённые количества, дефекты и исправления. Публикация следует действующей законной обязанности или разрешению, не новому отчёту каждого выпуска. Выписка проверяется на раскрытие и связывание; малые агрегаты могут идентифицировать человека. Произвольный единый размер ячейки не считается законной анонимностью.

Результаты не дают автоматического сокращения помощи. Неблагоприятное индивидуальное последствие требует законного решения и возможности исправления или пересмотра. Экспериментальные и теневые расчёты отделены от юридически значимых, пока компетентный акт не разрешит их использование. Проект не менял категорию реальной семьи.

\section{Исполнение, альтернативы регулирования и ресурсы}
Нужны ответственные за правовую спецификацию, источники, эталон, тесты, пересмотр и раскрытие. Работу могут вести существующие подразделения, но свободная мощность не доказана. Стоимость включает поля, снимки, защищённое хранение, эталон, обучение, доступ заявителя и повторный пересмотр. Наличие системы и открытого кода не означает нулевые затраты.

Первый шаг — подать в Агентство с копией Министерству материал 1 — нормативный проект; материал 2 — пояснительную записку; материал 3 — научное и правовое обоснование с вычислительным дополнением. Основная редакция — русская, с согласованными узбекской и английской версиями и сопроводительным обращением по статье 22. Имя и контакты взяты из предоставленного автором прежнего обращения; степени и экспертная регистрация не выдуманы. Статья 22 обосновывается воспроизводимым вкладом и технической демонстрацией, без права на назначение. Формальные требования обращения и участие в подготовке различаются. \cite{ref20}; \cite{ref7}.

\subsection{Альтернативы регулирования и ресурсы}
Действующий порядок (A0) сопоставляется с техническим исполнением законной основы (A1) и необходимым нормативным содержанием порученного акта (A2). A1 и A2 не взаимоисключающие варианты: проверки используют действующие полномочия, а правовое условие определяется компетентным актом либо сохраняется по точной применимой норме. Пять кратких пунктов материала 1 позволяют соразмерить форму. Это не измеренный рейтинг эффективности затрат; ручной резерв, новая отчётность и переход Президента по правам не предлагаются. \cite{ref1}; \cite{ref7}; \cite{ref2}; \cite{ref9}.

Пункт 2(b) ПП-5025 и статья 35 ЗРУ-682 предусматривают оценку регуляторного воздействия для мер, затрагивающих права и интересы граждан. Данное сравнение — вход для разработчика, не законченный официальный отчёт и не подтверждение всех согласований. Пояснение содержит действительных разработчиков, поручение, цель, механизмы, ожидаемый результат и сопутствующие изменения по пунктам 84–85 методики. \cite{ref19}; \cite{ref7}.

После сопоставления контролей, договоров и каналов оцениваются только непокрытые дополнительные задачи: правовое согласование, независимые случаи ответа, затронутые изменения интеграции/доступа и сопровождение. Часы отделены от дополнительных денег; перераспределённый труд не бесплатен. Агентство даёт объёмы и операционные доказательства, Министерство — законные средства. Исследование благосостояния отдельно. Цена модуля, экономия, свободный штат и изменение бюджета помощи не оценены.

\subsection{Последовательность внедрения и проверяемый учёт ресурсов}
Используются существующие правовые, системные и выпускные функции. До декабрьского внесения согласуются таблица решений и связь приёмочных доказательств со случаями; до первого применения с января — правило, обязательные согласования и проверки соответствия. Это зависимости, не оценка длительности или испытанная интеграция Агентства. Общенациональный ручной резерв не предписывается. Дополнительные договорные работы и средства требуют фактического определения; бюджетный контекст 2026 года не выделяет средства этому предложению на 2027 год.

Финансовая оценка должна учитывать объём работы, а не только поручать её существующему подразделению. Для каждой работы указываются единица, количество, дополнительные денежные затраты, часы работников, ответственный, доступная существующая мощность и источник финансирования. Отдельно выделяются согласование спецификации, интеграция, поддержка контроля, защищённое хранение, доступ, обучение, повторные проверки и пересмотр дел. Проверяемое учётное тождество:
\[
K_{\rm first}=K_{\rm setup}+K_{\rm annual},\qquad
K_{\rm annual}=K_{\rm fixed}+\sum_r Q_r k_r.
\]

Опубликованная основа оценки стоимости установлена: приложение 3 к постановлению № 687 распространяется на проекты Центра за счёт Фонда либо внебюджетных средств Агентства. Пункт 4 относится к сопровождению/совершенствованию систем, пункт 10 — к новым проектам; пункт 12 разграничивает условные параметры стоимости и фактический персонал. Материал 2 приводит формулу и условия применения. ПП-388 предусматривает программное бюджетирование, но не определяет цену этого модуля либо выделенное ему финансирование. Уточнены недостающие сведения: утверждённый состав и вид работ, модули, длительность, мощность и финансирование. \cite{ref24}; \cite{ref25}.
Здесь $Q_r$ — годовой объём в указанной единице, $k_r$ — дополнительные сумы на единицу; постоянные и переменные затраты не дублируются. Часы и мощность персонала учитываются отдельно, включая перераспределённый труд. Это определения учёта, не рассчитанная смета: необходимые объёмы и ставки Агентства не предоставлены.

Случаи выполняются в разрешённой среде выпуска, не дополнительным ручным этапом каждой заявки. Повторно используются доказанно равноценные проверки и обязательные внешние заключения. Изменение времени обработки требует реальной нагрузки и применимого испытания мощности; входы не получены. Законные сроки сохраняются. Равенство/земельные поправки и репрезентативное исследование имеют отдельные правовые, бюджетные и ресурсные требования; искусственные строки не определяют их цену.

\subsection{Необходимость нормативного содержания: отдельный тест не заменяет критерий}
PF-193, пункты 13–14, прямо поручает нормативный проект новой многомерной системы. Необходимость относится к законному содержанию оценки 2027 года, а не ко всем авторским техническим деталям. Материал 1 поэтому содержит пять кратких альтернативных пунктов для основного акта; они заменяют, а не дополняют полные семь пунктов и приложение. Краткая форма закрепляет правовую основу, соответствие исполнения, основания индивидуального решения, компетентное изменение и существующий порядок исполнения. Внутренние паспорта, случаи и журналы сохраняют техническую роль. \cite{ref1}.

По статьям 3, 14, 17, 18 и 37–40 ЗРУ-682 правовое условие определяется компетентным нормативным инструментом с применимым оформлением и публикацией. Ведомственный акт требует регистрации в предусмотренном случае. Правила № 3686, пункты 94–95, охватывают влияние на права и гарантии, межведомственную обязательность и действие вне системы органа; пункт 96(b) исключает организацию исполнения без новых норм и 96(v)–(j) устанавливает технические исключения. Поэтому техническое название не снимает нормативное содержание, а действующий технический документ без новых норм не требует превращения в акт. Пункт 20 требует необходимости и оценки альтернатив. Это поддерживает узкое включение, а не новое регулирование ради формы. \cite{ref7}; \cite{ref31}.

Правовое покрытие и техническое покрытие проверяются отдельно. Первое — точный источник, компетенция, сила, область, дата и обязательное последствие; второе — те же входы, даты, исключения, независимо обоснованный ответ и результат версии. Полное техническое покрытие снимает повторные тесты, но не доказывает отсутствующий источник критерия. Совпадение правового покрытия снимает дополнительную норму по этому вопросу. Если два юридически допустимых выбора дают разную категорию, только тест не выбирает законный вариант. Если существующая норма однозначно выбирает ответ, нужны её исполнение и, при необходимости, компетентное разъяснение, а не искусственная поправка.

Авторитетная область данных может исключить конкретный алгебраический вход; вывод тогда сужается. Программная валидация сама не доказывает правовой запрет семье сообщить факт. Статья 55 не разрешает менять акт толкованием. Пункт 48 приложения 1 № 35 прямо допускает временное изменение алгоритма при неисправностях, перерывах либо неполных электронных данных; это сохраняется и не трактуется как необходимость нового акта для каждого временного изменения. Предлагается запись его основания и области, а не расширение полномочия. Наличие такого механизма не доказывает определённость всей постоянной методики 2027 года. \cite{ref2}; \cite{ref7}.

Сокращение нормативного текста, отсылка к применимому правилу или законно уполномоченная ведомственная форма являются принятием по существу, если сохраняют обязательный результат, компетенцию, область и даты. Только техническое включение случаев не закрывает доказанно неурегулированное правовое условие. Интеграция использует действующие процессы; фактические часы, договорное покрытие и средства всё ещё не получены. Ресурсный аргумент способен изменить технический объём, но не легализует отсутствующий критерий. Не доказаны универсальная необходимость всех пяти пунктов, отсутствие равноценного действующего регулирования или уникальная оптимальность авторской формы. Это проверяемые пределы аргумента, а не обещание неизбежного принятия.

\subsection{Действующие контроли, соразмерность и законный выпуск}
Постановление № 35 уже закрепляет ответственность системы и источников, защиту, ограниченные полномочия при сбое и надзор. Пункт 6$^{1}$ постановления № 516 требует положительных технических/защитных заключений до внедрения; пункт 13 приложения 1 предусматривает совместные тесты подключения, пункт 44 — доступ к статистике запросов. Постановление № 93 вводит основу IT-аудита и функцию Министерства цифровых технологий. Контроли используются повторно. Успешный обмен сам по себе не доказывает правовой ответ на границе дохода/категории, а обязанность аудита не доказывает выполненный аудит расчётов Агентства. Такой реестр выпусков не получен. \cite{ref2}; \cite{ref26}; \cite{ref27}.

Дополнительный вклад — выполненные выводы и контрольные случаи с независимо обоснованным ответом, датой и результатом выпуска. Уже выполненная эквивалентная проверка используется повторно. Это снимает повторную техническую работу, но само по себе не доказывает нормативное покрытие условия: правовой источник проверяется отдельно. При подтверждении обоих покрытий дополнительная норма не требуется; при доказанно отсутствующем правовом условии необходима его компетентная нормативная фиксация в основном акте или законно предусмотренном инструменте. Общие обязанности и искусственные случаи не доказывают предельную пользу каждого пункта.

Пункты 7 и 23 исключают отдельную обязательную публичную отчётность. Утверждённое описание связывается на существующем ресурсе; доказательства сохраняются в существующих записях выпуска и надзора. Законные публикация, доступ заявителя и разрешённое раскрытие продолжаются. Любая разрешённая выписка проверяется на раскрытие, включая связывание малых агрегатов. Независимые контроли выполняются до внедрения или изменения, не вторым ручным расчётом каждой заявки. Не добавляется неоценённая общенациональная служба, но нулевая дополнительная трудоёмкость не заявляется.

Резерв Президента исключён из подаваемого нормативного проекта. Сохранение методики 2026 года после января, включение равенства дохода и освобождение от переходного долга создавали бы отдельные правовые и бюджетные решения; доказательства необходимости и доступности не получены. Необязательные исследовательские варианты равенства и земельного предела отдельно обозначены в материале 2 и не условия принятия. Пункт 26 теперь регулирует исправление и повторный тест, доказанное отделение незатронутых путей или другую соответствующую сборку только по действующему правилу. Нет полномочия ручного внесения, сохранения истёкшего правила или списания долга. Отсутствующая норма направляется компетентному разработчику/принявшему органу до первого применения; статья 55 не разрешает поправку толкованием. Исключение резерва не доказывает завершённость официального перехода. \cite{ref1}; \cite{ref2}; \cite{ref7}.

\subsection{Ответы на содержательные возражения}
Официальная спецификация баллов 2027 года для данного проекта не получена. Оценка нуждаемости и бюджета этой модели требует её получения; выполненный аудит относится к указанным опубликованным правилам и обосновывает предлагаемую проверочную обязанность. Если земельные входы не встречаются, это опровергало бы заявление о частоте, которого здесь нет; достоверные ограничения включаются в паспорт и тесты. Если приоритет тяжёлой ситуации уже действует, его точное правовое и программное основание разрешает вопрос, но не оправдывает замену специального пути общим интервалом.

Конфиденциальность обеспечивается законным внутренним расчётом и проверкой раскрытия, без передачи личных записей автору. Затраты требуют настоящего ресурсного расчёта и могут обосновывать этапность эмпирической работы; ответ «бесплатно» недопустим. Если действующая норма уже равноценно регулирует вопрос, разработчик сопоставляет её с предложенным пунктом и избегает повторения. Эти ответы сужают спор, не исключают отклонение и не обязывают принять предпочтительную политику.

Декабрьский срок — государственное поручение, не доказательство включения предложения. Различаются получение, встреча, назначение, принятие пункта, публикация для обсуждения и окончательное принятие. Успех подтверждается сопоставлением опубликованного акта с предложенными нормами. Из рабочей среды ничего не подано государству и принятие не состоялось.

\section{Установленные выводы и вопросы окончательной доработки}
Аудит охватывает выбранные формулы и условия, не всю систему. Нет микроданных, операционной программы, местных налоговых коэффициентов, окончательного округления и весов 2027 года. Земельная область сконструирована, переводы зависят от неразрешённого периода. Специальное толкование может разрешить тяжёлую ситуацию без изменения политики. Консолидированные страницы меняются: перед подачей проверяются повторно.

При этом официальный пример воспроизведён точно; выбранные условия признания и категории расходятся при \(P\) в заявленных условиях; ни одно простое прочтение общего интервала не сохраняет целиком специальный путь; буквальная площадь допускает отрицательность в заявленной области; вклад переводов зависит от периода. Налоговые и животноводческие результаты дают дополнительные эталоны. Число реально пострадавших или получивших выгоду не доказано.

Нормативный вклад — юридически значимый расчёт с определённой спецификацией, воспроизведением и пересмотром: единицы, периоды, исключения, версии, проверка границ до использования, сохранение неясности, основания решения и документированные контрольные расчёты. Это конкретный вклад в подготовку порученного акта. Эмпирическая эффективность и принятие требуют разрешённого анализа данных и компетентного государственного решения.

\renewcommand{\refname}{Источники}\begin{thebibliography}{99}
\bibitem{ref1} President of Uzbekistan. \textbf{PF-193}, 11 September 2026, paragraphs 13–14. \href{https://lex.uz/docs/-8472526}{Official consolidated text}.

\bibitem{ref2} Cabinet of Ministers of Uzbekistan. \textbf{Resolution 35}, 29 January 2026, Annex 1 paragraphs 18–30, 32–35, 37–42 and 48–49, and relevant administrative annexes. \href{https://lex.uz/docs/-8027513}{Official consolidated text}.

\bibitem{ref3} President of Uzbekistan. \textbf{PF-258}, 26 December 2025, paragraph 3 and Annex 1 notes (a)–(b). \href{https://lex.uz/docs/-7952180}{Official consolidated text}.

\bibitem{ref4} President of Uzbekistan. \textbf{PF-115}, 23 June 2026, minimum-remuneration provisions effective 1 September 2026. \href{https://lex.uz/docs/-8283656}{Official text}.

\bibitem{ref5} National Statistics Committee. \textbf{Minimum consumption expenditure for 2026}, 4 March 2026. \href{https://stat.uz/uz/matbuot-markazi/qo-mita-yangiliklar/67179-minimal-istemol-kharazhatlari-ijmati-t-risida-2}{Official announcement}.

\bibitem{ref6} National Statistics Committee. \textbf{Poverty rate, SIAT 3495}, 2021–2025 observations retrieved 30 September 2026. \href{https://siat.stat.uz/data/3495/}{Official data}.

\bibitem{ref7} Republic of Uzbekistan. \textbf{O‘RQ-682, On Normative Legal Acts}, 20 April 2021, Articles 3, 14, 17–18, 20–24, 32, 35, 37–40 and 55 and appended drafting methodology paragraphs 11, 43–51, 83–85. \href{https://lex.uz/docs/-5378966}{Official consolidated text}.

\bibitem{ref9} Republic of Uzbekistan. \textbf{O‘RQ-547, On Personal Data}, 2 July 2019, Articles 11, 16–18, 24, 27$^{1}$ and 28; special-data safeguards; Article 27$^{1}$ as amended by O‘RQ-1125 of 26 March 2026. \href{https://lex.uz/docs/-4396419}{Official consolidated text}.

\bibitem{ref10} Republic of Uzbekistan. \textbf{O‘RQ-457, On Administrative Procedures}, 8 January 2018, Articles 3, 53–54 and 59–61. \href{https://lex.uz/docs/-3492199}{Official consolidated text}.

\bibitem{ref11} World Bank. \textbf{Uzbekistan Poverty and Equity Assessment: Rising to Prosperity: Broadening Opportunities for All}, published 18 September 2026; printed pp. 40 and 78. \href{https://documents1.worldbank.org/curated/en/099091826081016988/pdf/P507037-a67dabb6-2302-46a0-be08-a951d35f90c3.pdf}{Report PDF}.

\bibitem{ref12} World Bank. \textbf{Global Insights on Social Registries: Coverage and Beyond}, 2025; printed pp. 13–15. \href{https://documents1.worldbank.org/curated/en/099071525165029379/pdf/P177331-a772f9c1-e340-4f76-aeb0-90fd4e7a3496.pdf}{Report PDF}.

\bibitem{ref13} Hadiwidjaja, Gracia; Williams, Asha; Giannozzi, Sara. \textbf{Improving Data Quality for an Effective Social Registry in Indonesia}, World Bank, 2022; printed pp. 22–23. \href{https://documents1.worldbank.org/curated/en/099515110132223290/pdf/P177341026be35061089cf0d835280eaf15.pdf}{Report PDF}.

\bibitem{ref14} Lindert, Kathy; Karippacheril, Tina George; Rodríguez Caillava, Inés; Nishikawa Chávez, Kenichi, editors. \textbf{Sourcebook on the Foundations of Social Protection Delivery Systems}, World Bank, 2020. DOI: 10.1596/978-1-4648-1577-5. \href{https://documents1.worldbank.org/curated/en/519831596182628993/pdf/Sourcebook-on-the-Foundations-of-Social-Protection-Delivery-Systems.pdf}{Book PDF}.

\bibitem{ref15} Grosh, Margaret; Leite, Phillippe; Wai-Poi, Matthew; Tesliuc, Emil, editors. \textbf{Revisiting Targeting in Social Assistance: A New Look at Old Dilemmas}, World Bank, 2022. \href{https://documents1.worldbank.org/curated/en/099300006162240907/pdf/P1648760df4c480270bd0b02288943f413e.pdf}{Book PDF}.

\bibitem{ref16} Alatas, Vivi; Banerjee, Abhijit; Hanna, Rema; Olken, Benjamin A.; Tobias, Julia. \textbf{Targeting the Poor: Evidence from a Field Experiment in Indonesia}. American Economic Review 102(4), 1206–1240, 2012. DOI: 10.1257/aer.102.4.1206. \href{https://economics.mit.edu/sites/default/files/publications/Targeting\%20Paper\%20AER\%20--\%20Final\%20Published.pdf}{Published author-hosted paper}.

\bibitem{ref17} 
Aiken, Emily; Bellue, Suzanne; Karlan, Dean; Udry, Chris; Blumenstock, Joshua E. \textbf{Machine learning and phone data can improve targeting of humanitarian aid}. Nature 603, 864–870, 2022. DOI: 10.1038/s41586-022-04484-9. \href{https://www.nature.com/articles/s41586-022-04484-9}{Publisher paper}.


\bibitem{ref18} 
Ministry of Justice. \textbf{Cabinet-submission drafting guide No. 2352}, 9 April 2012, current consolidated paragraphs 45–53. \href{https://lex.uz/docs/-1999242}{Official text}.


\bibitem{ref19} President of Uzbekistan. \textbf{PP-5025, On further improving regulatory impact assessment}, 15 March 2021, current paragraph 2(b). \href{https://lex.uz/docs/5331933}{Official consolidated Russian text}.

\bibitem{ref20} Republic of Uzbekistan. \textbf{On Appeals of Individuals and Legal Entities}, revised by O‘RQ-445, 11 September 2017, Articles 6 and 29. \href{https://lex.uz/docs/-3336169}{Official consolidated text}.

\bibitem{ref21} Barr, Earl T.; Harman, Mark; McMinn, Phil; Shahbaz, Muzammil; Yoo, Shin. \textbf{The Oracle Problem in Software Testing: A Survey}. IEEE Transactions on Software Engineering 41(5), 507–525, 2015; sections 2 and 5.1. DOI: 10.1109/TSE.2014.2372785. \href{https://discovery.ucl.ac.uk/id/eprint/1471263/1/06963470.pdf}{Author manuscript at UCL, retrieved 1 October 2026}.

\bibitem{ref22} Black, Paul; Guttman, Barbara; Okun, Vadim. \textbf{Guidelines on Minimum Standards for Developer Verification of Software}, NIST IR 8397, October 2021; sections 2.6–2.8, printed pp. 7–8. DOI: 10.6028/NIST.IR.8397. \href{https://nvlpubs.nist.gov/nistpubs/ir/2021/NIST.IR.8397.pdf}{Official report, retrieved 1 October 2026}.

\bibitem{ref23} OECD and European Commission Joint Research Centre. \textbf{Handbook on Constructing Composite Indicators: Methodology and User Guide}, OECD Publishing, 2008; section 1.7, printed pp. 34–35, and Step 7, printed pp. 117–131. DOI: 10.1787/9789264043466-en. \href{https://www.oecd.org/content/dam/oecd/en/publications/reports/2008/08/handbook-on-constructing-composite-indicators-methodology-and-user-guide_g1gh9301/9789264043466-en.pdf}{Official handbook, retrieved 1 October 2026}.


\bibitem{ref24} Cabinet of Ministers of Uzbekistan. \textbf{Resolution No. 687}, 28 December 2023, operative paragraph 3 and Annex 3 paragraphs 1, 3–4 and 7–12; consolidated version retrieved 1 October 2026. \href{https://lex.uz/docs/-6720230}{Official consolidated text}.

\bibitem{ref25} President of Uzbekistan. \textbf{PP-388}, 26 December 2025, paragraphs 19–21 and Annex 6; consolidated version retrieved 1 October 2026. \href{https://lex.uz/docs/-7959569}{Official consolidated text}.

\bibitem{ref26} Cabinet of Ministers of Uzbekistan. \textbf{Resolution No. 93}, 5 March 2026, operative paragraph 1 and annex paragraph 3: IT-audit framework and Ministry function. \href{https://lex.uz/docs/-8076376}{Official text checked 1 October 2026}.

\bibitem{ref27} Cabinet of Ministers of Uzbekistan. \textbf{Resolution No. 516}, 20 September 2022, operative paragraph 6$^{1}$ and annex 1 paragraphs 13 and 44; consolidated text checked 1 October 2026. \href{https://lex.uz/docs/-6200851}{Official text}.

\bibitem{ref28} Government portal. \textbf{Mahalla odamlar uchun ishonch, tayanch va imkoniyatlar markaziga aylanishi zarur}, 17 July 2026; paragraphs on Registry refusals, appeals and oversight. \href{https://gov.uz/oz/news/view/194551}{Official report checked 1 October 2026}.

\bibitem{ref29} Unified Interactive Public Services Portal. \textbf{Service 587: Ijtimoiy reyestrga kiritish yoki chiqarish uchun ariza berish}, service description and displayed status counters observed 1 October 2026. Counter accumulation period is not stated. \href{https://oldmy.gov.uz/oz/service/587}{Official service page}.

\bibitem{ref30} Urganch city administration. \textbf{O‘quv-amaliy seminar tashkil etildi}, 19 February 2026; local practical training report. \href{https://gov.uz/oz/urganchshahar/news/view/134892}{Official report checked 1 October 2026}.

\bibitem{ref31} Ministry of Justice of Uzbekistan. \textbf{Rules on preparation, adoption, legal examination and state registration of departmental normative legal acts}, registration No. 3686, 9 October 2025, paragraphs 4, 20 and 94–96; current version including amendments registered as No. 3943 on 17 September 2026, checked 1 October 2026. \href{https://lex.uz/docs/-7766987}{Official consolidated text}.

\end{thebibliography}

\appendix
\section{Повторение и границы аудита}
Вычисления повторяются без сети. Датированное извлечение правил и случаи — \nolinkurl{data/rules/legal_specification.json}; искусственный аддитивный набор — \nolinkurl{data/fixtures/additive_demo.json}. Это авторские извлечённые входы, не исходные записи семей и не полный сырой архив законодательства. JSON сохраняет SHA-256 спецификации и кода, команду, Python, UTC-время и детерминированность. Метаданные аудита и синтетического запуска также сохраняют путь действующего плана и его SHA-256. Случайной выборки нет, начальное случайное значение не нужно. Код — \nolinkurl{src/social_registry/}, задания — \nolinkurl{scripts/}, тесты — \nolinkurl{tests/}, входы — \nolinkurl{data/}, производные таблицы, графики, аудит, журналы и документы — \nolinkurl{results/}. Хеш относится к извлечению, не всей живой странице закона; перед подачей источник проверяется.

\begingroup\fontsize{11bp}{13.2bp}\selectfont
\begin{verbatim}
python scripts/run_analysis.py --self-test
python scripts/run_analysis.py --legal-audit \
  > results/audit/legal_rule_audit.json
python scripts/run_analysis.py --demo \
  > results/audit/synthetic_validation.json
python scripts/verify_results.py
python scripts/export_results.py
.venv/bin/python scripts/render_figures.py
.venv/bin/python scripts/build_documents.py --check-sources
\end{verbatim}
\endgroup
Движок использует стандартную библиотеку Python. Местная среда содержит matplotlib, Markdown, Beautiful Soup и pypdf для публикации. Полученные графики сохранены в results. Рукопись самостоятельна: таблицы, источники и монохромные векторные диаграммы встроены, последним соответствуют те же аналитические выражения. Компилятор редактора проверяет без дополнительных файлов. Необязательный терминальный сценарий компилирует поддерживаемый LaTeX и больше не восстанавливает его из Markdown; проверочный режим ничего не записывает. Прежний PDF перед подачей повторно создаётся или экспортируется из изменённого исходника.

Внешний ввод поддерживает только ограниченную аддитивную схему и не выполнялся на входе Агентства. Полные переводы, состав семьи, смешанное поголовье, права на программы, извлечение, защищённая инфраструктура и алгоритм 2027 года не реализованы. Административные агрегаты требуют отдельной проверки раскрытия.

\section{Сопроводительные документы}
Комплект организован как три пронумерованных материала: 1 — «Проект нормативных положений о проверке расчётов многомерной оценки семей»; 2 — «Пояснительная записка»; 3 — «Научное и правовое обоснование», включая электронное вычислительное дополнение. Сопроводительное обращение содержит просьбу об участии в подготовке по статье 22. Основной язык подачи — русский; узбекская и английская редакции являются согласованными переводами. Статья 6 закона об обращениях допускает подачу на других языках; статья 21 ЗРУ-682 регулирует подготовку официального проекта на государственном языке. Каждый основной материал содержит нормы или доказательства, необходимые для чтения без доступа к репозиторию. Вычислительное дополнение предоставляет план, датированное извлечение правил, контрольные случаи, реализацию, независимую арифметическую проверку и результаты выполнения. Официальный разработчик отвечает за окончательный правовой текст, ресурсы и процедуру принятия.

\end{document}
